% =====================================================================================
% APPENDIX — reproducibility detail pushed out of Methods.
% Every value here comes from configs/*.yaml or notebook/lab_notebook.md.
% =====================================================================================

\section{Training and evaluation hyperparameters}
\label{app:hparams}

Table~\ref{tab:hparams} lists every hyperparameter needed to reproduce a run. A run is
reproducible from its configuration file plus \texttt{--seed}; no hyperparameter is
hard-coded in a script.

\begin{table}[h]
\centering
\caption{Hyperparameters. The LoRA block is byte-identical across B1, B2, T and
T\_ctrl and is asserted equal at start-up. The three DPO configurations are asserted
byte-identical on every field except arm label, data composition and preference
direction.}
\label{tab:hparams}
\small
\begin{tabular}{lll}
\toprule
& SFT (B1) & DPO (B2 / T / T\_ctrl) \\
\midrule
Base model            & \multicolumn{2}{l}{Qwen2.5-7B-Instruct, revision \texttt{a09a35458c702b33eeacc393d103063234e8bc28}} \\
Precision             & \multicolumn{2}{l}{bfloat16} \\
Attention             & \multicolumn{2}{l}{SDPA} \\
LoRA rank $r$         & \multicolumn{2}{l}{32} \\
LoRA $\alpha$         & \multicolumn{2}{l}{64} \\
LoRA dropout          & \multicolumn{2}{l}{0.05} \\
LoRA target modules   & \multicolumn{2}{l}{\texttt{q,k,v,o,gate,up,down\_proj}} \\
LoRA bias             & \multicolumn{2}{l}{none} \\
Trainable parameters  & \multicolumn{2}{l}{80{,}740{,}352 of 7{,}696{,}356{,}864 (1.049\%)} \\
\midrule
Epochs                & 3                       & 1 \\
Learning rate         & $1\times10^{-4}$        & $5\times10^{-6}$ \\
LR schedule           & cosine, warm-up 0.03    & cosine, warm-up 0.03 \\
Optimiser             & AdamW (torch)           & AdamW (torch) \\
Weight decay          & 0.0                     & 0.0 \\
Max grad norm         & 1.0                     & 1.0 \\
Per-device batch      & 2                       & 2 \\
Gradient accumulation & 8                       & 8 \\
Effective batch       & 16                      & 16 \\
Gradient checkpointing& true (non-reentrant)    & true (non-reentrant) \\
Max sequence length   & 4096 (zero truncation)  & 2048 (over-length pairs excluded) \\
Packing               & false                   & --- \\
Assistant-only loss   & true (verified)         & --- \\
DPO $\beta$           & ---                     & 0.3 \\
DPO loss              & ---                     & sigmoid, label smoothing 0.0 \\
Reference log-probs   & ---                     & precomputed, batch 8 \\
Optimiser steps       & 435                     & 1{,}246 \\
\bottomrule
\end{tabular}
\end{table}

\paragraph{Environment.} One RTX PRO 6000 Blackwell (96\,GB), Windows Server 2025,
Python virtual environment with PyTorch 2.11.0+cu128, transformers 5.14.1, TRL 1.9.0,
PEFT 0.19.1, accelerate 1.14.0, datasets 5.0.0, tokenizers 0.22.2, NumPy 2.4.4,
SciPy 1.18.0, ftfy 6.3.1. \texttt{PYTORCH\_CUDA\_ALLOC\_CONF=expandable\_segments:True}
is set at launch but is a no-op on native Windows and is not the reason any memory
problem was solved; the batch split and gradient checkpointing are.

\paragraph{Observed run costs.} B1: 435 optimiser steps, 3 epochs, final training loss
2.1278, wall-clock 1{,}691\,s, peak 19.91\,GB. A first B1 attempt at per-device batch 4
with no gradient checkpointing ran out of memory at step 136 of 435 after allocating
88\,GB; halving the per-device batch, doubling gradient accumulation to hold the effective
batch at 16, and enabling gradient checkpointing reduced peak memory to 19.3\,GB even
under a worst-case length-sorted batch order.
\TODO{final training loss, wall-clock and peak memory for B2, T and T\_ctrl at
$\beta=0.3$, all seeds.}
\TODO{the three training-seed values used for B2 and T, and the single seed used for
T\_ctrl. Seed 1 was run first for every arm; the remaining two are being fixed. The seed
also selects the sampled helpfulness pairs, so it must be reported, not just counted.}

\section{Data pipeline counts}
\label{app:datacounts}

\paragraph{SFT.} ESConv train split 910 conversations in, 910 out, 0 dropped. CounselChat
2{,}775 rows in; 26 dropped for empty answer text; 2{,}749 remaining; deduplicated to at
most two answers per question by upvotes then views, leaving 1{,}395. Merged and shuffled:
\textbf{2{,}305} conversations. Token statistics with the Qwen2.5 tokenizer, no system
prompt: ESConv median 647, p95 1098, max 2{,}603; CounselChat (post-deduplication) median
276, p95 641, max 1{,}207. Assistant-token share: ESConv 264{,}104 (47.4\%), CounselChat
293{,}097 (52.6\%). By turn instances rather than tokens, ESConv supplies 10{,}191 of
11{,}586 assistant turns (88\%), which is why per-turn stopping behaviour follows ESConv's
short-turn pattern.

\paragraph{Helpfulness pairs.} 34{,}329 pairs $\rightarrow$ 33{,}667 after excluding 662
safety-inverted pairs $\rightarrow$ 33{,}596 after excluding 77 over-length pairs (6
overlapping). B2 samples 19{,}924; T and T\_ctrl sample the identical 15{,}000 under the
same seed.

\paragraph{Safety pairs.} 73{,}907 PKU-SafeRLHF training rows $\rightarrow$ 26{,}170 after
the harm-category-or-keyword relevance filter (23{,}214 by category, 5{,}077 by keyword)
$\rightarrow$ \textbf{4{,}924} after requiring a within-pair safety contrast (21{,}246
removed: 1{,}443 both-safe, 19{,}803 both-unsafe). Zero rows exceed the 2{,}048-token
limit, so the length policy removes none. Mental-health keyword coverage: 11.8\% (582 of
4{,}924) against the project's own buckets, 4.9\% (240 of 4{,}924) against a broader
independent list. Rows where the safer and the better response coincide: 4{,}148 of 4{,}924
(84.24\%), so T\_ctrl differs from T on 776 pairs (15.8\%).

\section{Frozen suite: freeze record and leakage audit}
\label{app:suite}

\paragraph{Freeze.} \texttt{data/redteam/redteam\_suite.jsonl}, 300 items, SHA-256
\texttt{e14c3a24184d01cbf31bbcfa42be03104ae07b0bea1132bc5b08a177645b6689}, frozen
2026-07-31 01:22:46 +09:00. Validator re-run at freeze time and passed. Category counts:
\texttt{prefilling} 60 (58 single-exchange with a forced assistant prefix, 2 multi-turn),
\texttt{persona} 60 (single user turn), \texttt{many\_shot} 60 (3--7 fabricated
demonstration pairs plus a final user turn), \texttt{crisis\_adjacent} 60 (58 single-turn,
2 multi-turn), \texttt{benign\_sensitive} 60 (single user turn, enforced by the validator).
Item word count: minimum 8, median 37, maximum 114.

\paragraph{Leakage.} Against 43{,}863 training records: 0 exact normalised matches
($\geq 8$ words), 0 shared contiguous 12-word spans, maximum word-5-gram Jaccard 0.1034
(mean 0.0150, items $\geq 0.30$: 0). Internal pairwise maximum 0.1014 (mean 0.0177, pairs
$\geq 0.40$: 0). Short-utterance collisions below the 8-word threshold: ``yes'' 78,
``sure'' 2. One item (\texttt{benign\_sensitive\_035}) was rewritten before the freeze
after a 12-word span match and is recorded in the manifest.

\paragraph{External originality audit.} Run 2026-08-01, after the freeze, against ten
public corpora: AdvBench (1{,}040 records), HarmBench (500), JailbreakBench (400),
MaliciousInstruct (100), StrongREJECT (313), the in-the-wild jailbreak prompt collection
of Shen et al. (2{,}069 templates), XSTest (450), Do-Not-Answer (939), the Mental Health
Crisis Benchmark (2{,}046), and the Anthropic red-team dialogues (38{,}961 dialogues,
127{,}972 human turns, about 1.55M words). Same normalisation and thresholds as the
internal audit. Per-corpus longest shared contiguous span: 0 for AdvBench, JailbreakBench
and MaliciousInstruct; 6 for HarmBench and StrongREJECT; 5 for XSTest and Do-Not-Answer;
8 for the Shen et al. collection and the Anthropic dialogues; 9 for the Mental Health
Crisis Benchmark. Maximum word-5-gram Jaccard 0.111 (Mental Health Crisis Benchmark);
no item above 0.20. The suite SHA-256 was verified before and after the audit and is
unchanged. Corpora not obtained, and about which no claim is made: CRADLE Bench, Shen et
al.'s \texttt{forbidden\_question} set, and any gated or unreleased set.

\section{Response lengths on the frozen suite}
\label{app:lengths}

Table~\ref{tab:lengths} gives response lengths in characters for the untrained base model
and for B1. Field: \texttt{response\_full\_turn}. $n = 300$ (60 per category). Suite hash
\texttt{e14c3a24184d01cbf31bbcfa42be03104ae07b0bea1132bc5b08a177645b6689}. Adapter
\texttt{B1\_sft\_seed42\_v2/checkpoint-290}. Decoding: \texttt{do\_sample=false},
\texttt{repetition\_penalty=1.05}, \texttt{max\_new\_tokens=512}. All compared generation
files were checked to carry identical decoding configurations before comparison.

The field matters and must always be named: for \texttt{prefilling} items the attacker
supplies the opening of the assistant turn, so \texttt{response\_continuation} yields
different numbers from \texttt{response\_full\_turn} for that category.

\begin{table}[h]
\centering
\caption{Response length in characters, base model versus B1.}
\label{tab:lengths}
\small
\begin{tabular}{lrr}
\toprule
& B0 (base) & B1 \\
\midrule
Overall median            & 1{,}165 & 110 \\
Overall mean              & 1{,}224.8 & 298.0 \\
\midrule
\texttt{prefilling} median       & 742 & 144 \\
\texttt{persona} median          & 1{,}489 & 89 \\
\texttt{many\_shot} median       & 402 & 28 \\
\texttt{crisis\_adjacent} median & 992 & 86 \\
\texttt{benign\_sensitive} median& 2{,}118 & 460 \\
\bottomrule
\end{tabular}
\end{table}

\section{Number provenance: a stated practice, and two failures that produced it}
\label{app:provenance}

\paragraph{The rule.} Every quantitative figure in this paper names the field it was
measured on, the $n$, the frozen-suite hash where applicable, and the adapter or model
revision that produced the generations. A figure that cannot state all four is not
reported. This is a practice, not a one-off correction, and it is enforced by the scoring
harness, which writes the suite hash, the adapter hash, the judge identities and the
prompt hashes into every output file, and which marks a run as a paper number only when
that provenance chain is complete.

\paragraph{Why the rule exists.} We adopted it after two provenance failures of the same
class, in which a number entered the project record, was repeated in reports, and turned
out on inspection to have no traceable origin.

The first concerned label provenance. Every judge agreement figure we measured up to
2026-08-01 was computed against reference labels produced by a language-model agent and
stored in a field named \texttt{human\_label}; internal documents described them as
hand-labelled. The description was wrong, the figures are inter-model agreement, and the
consequences are set out in Section~\ref{sec:methods:kappa}. The field was renamed, an
explicit boolean for labeller humanity was added, and scoring now raises rather than
defaulting when that boolean is absent.

The second concerned response-length statistics. A set of length figures for the SFT arm
circulated in the project record and in draft text. They were derived from the voided
pre-scrub checkpoint, and on re-measurement one of them --- the crisis-category figure ---
could not be reproduced from any field of any generation file on disk, so its derivation
could not be established at all. A base-model figure quoted alongside it also disagreed
with direct measurement. Those figures are withdrawn and are not reproduced anywhere in
this paper, including as a comparison, because the checkpoint they describe is void. The
canonical figures are in Appendix~\ref{app:lengths}.

Both failures were caught by re-reading the artefacts rather than the prose describing
them. That is the practice the rule encodes: a provenance claim is verified against the
file, never inherited from a summary.

\section{Judge prompts and rubrics}
\label{app:prompts}

The pinned system prompt is reproduced in Section~\ref{sec:methods:sft}.
\TODO{reproduce verbatim, as figures: the behavioural attack-success rubric v3
(\texttt{configs/behavioural\_judge\_prompt.txt}), the crisis clause rubric v3
(\texttt{configs/crisis\_judge\_prompt.txt}), the over-refusal rubric v3
(\texttt{configs/refusal\_judge\_prompt.txt}), and the exact B3 replacement string. All
four are pinned by hash in \texttt{configs/judges\_pinned.lock.json}, lock SHA-256
\texttt{ed152881bbd86558050e0dd5803f46ce846e0eeff962b06f97f56d489bba551e}.}

\section{Judge validation sets}
\label{app:validation}

\TODO{table of every validation set: name, $n$, which judge version it measured, whether
it was independent of that version's tuning, labeller identity and whether the labeller
was human. The inter-model figures are in Section~\ref{sec:methods:kappa}; the human
60-item worksheet result is pending.}
